\section*{Discussion}
Cell-cell interactions play a crucial role in various biological processes, including tissue development, homeostasis, and disease. These interactions are essential for the proper functioning of multicellular organisms, as they facilitate communication and coordination among different cell types. Single-cell RNA-seq approaches have allowed us to comprehend the heterogeneity of cell types in a complex tissue. However, traditional single-cell RNA-seq data lacks spatial context, making it challenging to uncover the spatial organization and intercellular relationships that underpin complex biological systems. Spatial transcriptomics is a burgeoning field that addresses this limitation by providing spatially resolved gene expression data, enabling researchers to explore cells in a spatial context and bridging the gap between molecular and morphological observations.\par
 
In this regard, our newly developed computational method, Renoir, offers a powerful tool for investigating cellular interactions within their native spatial context, from development to disease. Renoir enables the spatial mapping of ligand-target activities through the quantification of a neighborhood activity score, infers spatial communication domains by grouping spots based on the similarity of ligand-target activities and uncovers ligand-target interactions restricted to each communication domain. Renoir further enables the identification of the major ligands active in each domain and the spatial mapping of pathway activities based on known or denovo genesets. Renoir can work on datasets generated using spot-resolution spatial technologies as well as single-cell resolution spatial technologies. \par

Our comprehensive benchmarking of Renoir on a variety of datasets simulated from real datasets across tissue types demonstrate its superiority over existing CCC inference methods in quantifying the ligand-target activities across spatial locations. Specifically, by accounting for the receptor expression in the target cell type, Renoir is capable of reducing false positive interactions which are abundantly found by other methods. Using DLPFC spatial datasets with ground truth annotations of spatial domains, we further illustrated that the communication domains inferred by Renoir can better capture the underlying spatial structure of the tissue as compared to other methods. \par 

Our application of Renoir to three distinct 10X Visium spatial transcriptomics datasets including adult mouse brain, human fetal liver, and human triple negative breast cancer (TNBC) demonstrate Renoir's versatility in characterizing ligand-target interactions across tissue types. For each dataset, Renoir identified distinct spatial communication domains which harbored distinct ligand-target activities pertaining to the major cell types abundant in the domain and these were also consistent with known literature. In the adult mouse brain, Renoir characterized the interactions of different regional astrocyte subtypes with other co-localized excitatory and inhibitory neurons. For the triple negative breast cancer, Renoir uncovered spatial communication domains and sub-domains harboring distinct chemokine interactions between cancer, stromal and immune cells which were associated with tumor proliferation and metastasis. Renoir further identified signaling pathways that were active in cancer (WNT signaling) or immune/stromal cell type (IL-6/JAK/STAT3 signaling) dominated domains of the TNBC tissue. For developing fetal liver, Renoir uncovered an interaction between hepatocytes and macrophages mediated via the ligand PLG. The application of Renoir on a Nanostring CosMx hepatocellular carcinoma dataset further shows Renoir's utility in handling single-cell resolution large-scale spatial datasets. From the hepatocellular carcinoma dataset, Renoir identified several onco-fetal interactions involving endothelial cells, macrophages and CAFs and uncovered the role of these onco-fetal cell types in activating prostaglandin signaling in the bipotent hepatocytes. \par

By integrating single-cell RNA-seq and spatial transcriptomics data, Renoir enables a more comprehensive understanding of the dynamic relationships between different spatially co-localized cell types and their roles in various biological processes. This approach allows for the identification of key players in cellular communication and provides insights into the mechanisms governing tissue organization, functionality, and response to disease. The several visualization modules provided by Renoir further enable the user to explore the spatial map of ligand-target and pathway activities of their interest. Ultimately, Renoir has the potential to advance our understanding of complex biological systems and contribute to the development of novel therapeutic strategies targeting cellular interactions.